\documentclass[10pt,journal]{IEEEtran}
\usepackage{microtype}
\usepackage{amsmath,amssymb,mathtools}
\usepackage{booktabs,multirow,array}
\usepackage{graphicx}
\usepackage{cite}
\usepackage{xcolor}
\usepackage{url}
\emergencystretch=1em
% Retain IEEEtran's journal baseline; page length is controlled by prose and
% float placement without compressed line spacing.
\setlength{\textfloatsep}{3pt plus 1pt minus 1pt}
\setlength{\floatsep}{2pt plus 1pt minus 1pt}
\setlength{\intextsep}{3pt plus 1pt minus 1pt}
\setlength{\abovedisplayskip}{2pt plus 1pt minus 1pt}
\setlength{\belowdisplayskip}{2pt plus 1pt minus 1pt}
\renewcommand{\arraystretch}{0.98}
\newcommand{\RiskSafe}{risk-constrained verifier}
\newcommand{\TailRisk}{tail-risk counterfactual}
\newcommand{\PaymentSafe}{payment-certified verifier}
\newcommand{\dt}{\Delta t}

\title{Security-Aware Workload-State Counterfactual Verification and\\
Nodal Net-Value Settlement for Spatially Coupled AI Data Centers}
\author{Anonymous Authors}

\begin{document}
\maketitle

\begin{abstract}
Artificial-intelligence (AI) data centers can provide demand response (DR), but meter decreases do not show pre-event workload commitment or positive network value. Delayed service and cross-site rebound make workload-feasibility and signed-value checks essential. This paper develops a declaration-indexed verifier for spatially coupled data centers. Submit-time scheduler records, including release, declared runtime, graphics processing unit (GPU) request, class, and region, define a cumulative workload-state model, a convex risk contract, and a contiguous runtime witness on one ledger. An active-window energy map includes carry-in jobs, while execution telemetry remains outside decision-time fitting. The risk fit combines total exposure with conditional value-at-risk (CVaR) under a validation-frozen mean-square-error (MSE) bound; exact job-start enumeration audits its finite executable bridge. In a 54-day mechanism-isolation panel, the joint fit reaches normalized root-mean-square error (nRMSE) 0.902, false credit 0.780 MWh/day, under-credit 26.506 MWh/day, and a precision--recall harmonic-mean $F_1$ score of 0.536. The total-budget-only fit has lower nRMSE, false credit, and under-credit, higher $F_1$ (0.549 versus 0.536), and lower false-credit CVaR. Thus, these locked results do not establish an empirical benefit from the tail term. At the slot-62 information gate, nRMSE is 1.366 and $F_1$ is 0.080. The job-level bridge exactly audits 7,306,622 feasible starts, while its flexible-load relative-$L_2$ residual is 1.581. The IEEE Reliability Test System 24-bus (RTS-24) result is a declaration-only single-line N--1 network replay; payment activation requires a closed meter, and observational workload traces do not identify a utility intervention effect.
\end{abstract}

\noindent\textbf{Index Terms---}AI data centers, counterfactual verification,
demand response, security-constrained dispatch, nodal settlement.

\section{Introduction}

The growth of large-model training and inference is converting data centers
from nearly static industrial loads into large, software-defined grid
resources. Interactive inference is latency constrained, elastic inference can
be delayed briefly, and batch GPU work can be deferred for hours or migrated
between regions. These capabilities enable
demand response (DR), but they also
invalidate the usual interpretation of a meter decrease. A reduction at one
facility can be followed by delayed consumption or simultaneous rebound at
another node. Paying the positive local decrease can therefore reward a
redistribution that creates little or negative network value.
Public traces already expose multi-week bursty workload arrivals and tens of
thousands of scheduler jobs, while facility telemetry reports GPU energy at a
different measurement boundary \cite{wang2024burstgpt,samsi2021supercloud}.

Conventional DR measurement and verification estimates the unobserved
no-event demand from historical meter data. Arithmetic historical baselines
remain common in market practice under published market guidance
\cite{caiso2017baseline}, despite known bias and manipulation incentives
\cite{wang2022baseline}. Regression-based facility models quantify
baseline uncertainty but do not encode a computing workload ledger
\cite{wang2022baseline}. More broadly, DR
research establishes the economic role of responsive load, yet a meter-only
estimator cannot determine whether queued jobs could actually have produced
its predicted trajectory.

Data-center scheduling studies provide this physical state: temporal deferral
and geographical balancing reduce energy cost and peak demand
\cite{cao2022flexibility}. Data-driven flexibility and
security-constrained coordination optimize spatial--temporal dispatch
\cite{cao2022flexibility,cao2024nonwire}. Balancing-market and load-shifting
studies provide complementary dispatch contexts. Virtual-link clearing and
space--time remuneration model geographical and temporal transfers
\cite{zhang2020virtuallinks,zhang2022remunerating}. Receding-horizon frequency
regulation retains multi-site state and future arrivals
\cite{zhang2023receding}. These
methods optimize operator-controlled schedules; verification of an independently
submitted DR counterfactual against an auditable workload ledger remains open.
Coupled-regulation models likewise provide no provenance-bound certificate
\cite{chen2021idccoupling}. Whole-facility measurements
also require a bottom-up GPU-to-facility conversion; calibration, benchmark
scaling, and network placement are therefore kept as separate evidence layers.

The second gap is electrical. Geographic workload control is often valued with
energy prices or gross megawatts, although nodal marginal value changes when
congestion or a contingency binds. Direct-current (DC) power flow and linear distribution
factors provide transparent sensitivities \cite{stott2009dc,zimmerman2011matpower};
N--1 dispatch imposes credible post-contingency limits through line-outage
distribution factors (LODFs) \cite{tejada2018lodf}. No data-center DR verifier
jointly certifies workload feasibility, signed spatial response, and
security-aware network value. Alternating-current (AC) feasibility is a separate
nonlinear check. Public traces support observational
replay but not a causal utility event, so measured execution, counterfactual
construction, and network valuation must remain separate.

This paper makes three power-system contributions. 1) It defines a typed
submit--execute--settle state in which a runtime-complete job witness is formed
from release, declared runtime, GPU request, and regional assignment. An
active-window energy map gives the exact declaration energy intersecting a
locked operating window, including carry-in jobs. The same indexed digest is
rechecked at regional aggregation and reused by the RTS-24 N--1 value replay.
2) It formulates a workload-state verifier with a causal information boundary,
a validation-selected high-5-of-10 Pareto projection, and an exact convex
total-plus-daily-CVaR contract. A predeclared 10\% MSE neighborhood and a
triangle-inequality certificate prevent risk reduction from being purchased by
an unreported accuracy change; the full simplex is solved numerically and
audited with Karush--Kuhn--Tucker (KKT) and primal residuals. The aggregate contract is then mapped to
a fixed linear dual-price surrogate and an exact finite contiguous-block policy;
the resulting target residual is reported rather than assumed away. 3) It gives
the workload response a signed nodal value under complete finite N--1
contingencies and separates declaration feasibility, network-value replay, and
meter-gated payment activation. This separation prevents a declaration
trajectory from being interpreted as a closed-meter payment. The resulting
coupling invariant carries one declared workload state into physical
verification and security-aware valuation; its LP, dispatch, and allocation
primitives retain established power-system semantics \cite{li2024computation}.
Because public traces do not identify utility geography, spatial conclusions
are conditional on the declared feature-stratified trace scenario and the
complete permutation/penetration panel. The certificate assumes trusted
scheduler provenance: a digest detects record changes but cannot authenticate
semantically false fields.

The remainder of the paper is organized as follows. Section~\ref{sec:model}
defines the spatial workload and network settlement problem. Section~\ref{sec:method}
presents the verifier, security-aware value rule, and allocation mechanism.
Section~\ref{sec:design} specifies the data, baselines, and locked protocol.
Section~\ref{sec:results} reports the core counterfactual, execution-replay,
job-level, security, payment, continuous-horizon, and provenance experiments.
Section~\ref{sec:conclusion} concludes with the deployment boundary and future
validation requirements.

\section{System Model and Problem Formulation}\label{sec:model}

\subsection{Spatial workload state}

Let $\mathcal S$, $\mathcal K$, $\mathcal D$, and $\mathcal T$ denote source
regions, workload classes, destination data centers, and 15-minute intervals;
the symbol $h$ is reserved for a calendar day in statistical panels.
The submit-time declaration ledger supplies arrival energy $a_{skt}$ for source
$s$, class $k$, and release interval $t$. The decision $x_{skd\tau}\geq0$ is the energy
served at destination $d$ in interval $\tau$. Its cumulative state is
\begin{equation}
 y_{skt}=\sum_{\tau=0}^{t}\sum_{d\in\mathcal D}x_{skd\tau}.
 \label{eq:cumulative}
\end{equation}
With cumulative arrivals $A_{skt}=\sum_{\tau=0}^{t}a_{sk\tau}$ and deadline
$D_k$, release and deadline feasibility are
\begin{align}
 y_{skt}&\leq A_{skt}, \label{eq:release}\\
 y_{skt}&\geq A_{sk,t-D_k},\quad t\geq D_k. \label{eq:deadline}
\end{align}
The state recursion is $y_{skt}-y_{sk,t-1}=\sum_d x_{skdt}$ with
$y_{sk,-1}=0$, and terminal service equals the declared arrivals due in the
window; the release and deadline inequalities are (\ref{eq:release}) and
(\ref{eq:deadline}), and the site-capacity constraint is
$\sum_{s,k}x_{skdt}\leq\bar P_d\dt$. Facility demand is
\begin{equation}
 p_{dt}=P_d^{\rm fix}+\dt^{-1}\sum_{s,k}x_{skdt}.
 \label{eq:power}
\end{equation}
For rolling window end $T^+$ and final event-day release $T^0$, the terminal
state is not forced to consume unexpired future work. Instead,
\begin{equation}
y_{skT^+}=A_{sk,\max\{T^0,T^+-D_k\}},
\label{eq:rollingterminal}
\end{equation}
which completes every job due inside the window and every event-day arrival
while retaining later real arrivals. The window starts $D_{\max}+96$ intervals
before and ends $D_{\max}$ intervals after the event; omitted earlier work has
expired. Equations (\ref{eq:release})--(\ref{eq:power}) map an aggregate preemptive
envelope to the facility meter. This is valid for checkpointable batch work; a separate execution audit checks
contiguous intervals, runtime, GPU count, and native power.
For the counterfactual task-level check, let \(j\in\mathcal J_{\rm sub}\) index a
submitted scheduler record, \(r_j\) denote its submitted release, and let
\(d_j\) be the service-window endpoint. Slurm ``timelimit'' is an allocation run-time
declaration; it is converted to \(R_j=\lceil\mathrm{timelimit}_j/\dt\rceil\)
slots and combined with a precommitted queue allowance \(B=96\) slots, so
\(d_j=r_j+R_j+B\). The unlimited Slurm sentinel is mapped to the precommitted
horizon \(H_{\infty}=128\) slots before adding \(B\); finite declarations are
retained exactly. The $0.001$-MW-per-GPU nameplate is fixed before the split.
The training-calibrated entitlement \(\widehat E_j\) is retained only for the
aggregate Exp19 audit; the executable witness uses the declared runtime. Let
\(\bar p_{\rm GPU}=0.001\) MW/GPU be the committed nameplate, $s_j$ the
submit-time scenario region, and \(g_j^{\rm req}\) the requested GPU count.
Define the runtime-complete entitlement
\(E_j^{\rm run}=R_jg_j^{\rm req}\bar p_{\rm GPU}\dt\). The indexed service
variable \(u_{j\tau}\) is the energy assigned to job \(j\) in interval \(\tau\),
and the executable witness requires one contiguous fixed-rate block of exactly
\(R_j\) slots. Let \(w_{\rm batch}\) be the predeclared waiting-cost coefficient per
MWh-slot, and \(r^{\rm DR}\) the event service tariff. Define
\(c_j=g_j^{\rm req}\bar p_{\rm GPU}\dt\) and
\(\Theta_j=\{r_j,\ldots,d_j-R_j\}\). For each admissible start \(\theta\),
a binary \(z_{j\theta}\) selects the block and
\(u^{j\theta}_{\tau}=c_j\mathbf1\{\theta\leq\tau<\theta+R_j\}\). The exact
start-time mixed-integer linear program (MILP) is
\begin{equation}
\begin{aligned}
 &\sum_{\theta\in\Theta_j}z_{j\theta}=1,\qquad z_{j\theta}\in\{0,1\},\\
 &u_{j\tau}=\sum_{\theta\in\Theta_j}u^{j\theta}_{\tau}z_{j\theta},\qquad
   \sum_{j:s_j=d}u_{j\tau}\leq\bar P_d\dt,\\
 &\min_z\ \sum_{j,\theta\in\Theta_j}z_{j\theta}
   \sum_{\tau}\left[w_{\rm batch}(\tau-r_j)+r^{\rm DR}\mathbf1\{\tau\in\mathcal E\}\right]u^{j\theta}_{\tau}.
\end{aligned}
\label{eq:jobindexed}
\end{equation}
Every selected block enforces runtime completion, contiguity, and the GPU
nameplate; positive regional slack makes the separable enumeration exact for
the displayed MILP. For delay
$\delta=\theta-r_j$, the waiting part of the objective is
$w_{\rm batch}c_j\{R_j\delta+R_j(R_j-1)/2\}$, so a long job is charged for every
delayed service interval rather than only for its start time. The runtime-complete
ledger enumerates every admissible start for all 75,326 records and stores the
earliest-start baseline and tariff response; a separate capacity case activates
a binding regional row. The measured
\([t_j^{\rm start},t_j^{\rm end}]\) interval is opened only for independent
post-event scoring after the declaration witness is fixed: 71,128 execution
matches never constrain \(z\), \(u\), or \(d_j\). Thus \(\widehat E_j\) is
calibration-only, \(E_j^{\rm run}\) is the nameplate-capacity commitment, and
the calibrated central energy is carried on the same blocks without being
presented as measured consumption. The risk target is linked through a fixed
regional slot price. Because the statistical and nameplate ledgers have
different scales, define the locked window set
$\mathcal W_{\rm lock}=\bigcup_{h\in\mathcal H_{\rm lock}}
\{96h,\ldots,96h+95\}$. The declaration energy active in this set is
\begin{equation}
 E_{\rm act}(\mathcal W_{\rm lock})=
 \sum_{j\in\mathcal J_{\rm sub}}E_j^{\rm dec}
 \frac{|[r_j,r_j+R_j)\cap\mathcal W_{\rm lock}|}{R_j},
 \label{eq:activewindow}
\end{equation}
where $E_j^{\rm dec}$ is either the nameplate upper entitlement or the
calibrated central entitlement. This overlap integral in (\ref{eq:activewindow}) includes carry-in jobs
and uses only submit-time fields. Let
$\widetilde p^{\rm risk}_{dht}=P_d^{\rm fix}+\gamma_+[\bar p^{\rm risk}_{dht}-P_d^{\rm fix}]_+$,
with $\gamma_+=E_{\rm act}(\mathcal W_{\rm lock})/\sum_{d,h,t}
[\bar p^{\rm risk}_{dht}-P_d^{\rm fix}]_+\dt$. For locked day $h$, region
$d$, and interval $t$, define the fixed linear dual-price surrogate
$\pi_{dht}=r^{\rm DR}[\widetilde p^{\rm risk}_{dht}-P_d^{\rm fix}]_+ /
\widetilde\Pi_{\max}$, where $\widetilde\Pi_{\max}=\max_{d,h,t}
[\widetilde p^{\rm risk}_{dht}-P_d^{\rm fix}]_+$. Re-enumerating every
declaration-feasible block with the waiting, event-tariff, and $\pi_{dht}$
terms produces the exact minimizer of this finite linear-price objective. The
optimality statement concerns the declaration-feasible policy; equality to the
aggregate target is a separate primal residual. The bridge therefore records,
for every locked day,
\begin{equation}
 \begin{aligned}
 e^{\rm br}_{dht}&=p^{\rm exec}_{dht}-\widetilde p^{\rm risk}_{dht},\\
 \mathrm{nRMSE}_{\rm flex}&=
 \frac{\|e^{\rm br}\|_2}{\max\{\|\widetilde p^{\rm risk}-P^{\rm fix}\|_2,\epsilon\}}.
 \end{aligned}
 \label{eq:bridge_residual}
\end{equation}
along with the total-load normalization and the flexible relative-L2
denominator. An independent enumeration checks all 7,306,622 candidate starts,
verifies the stored minimizer, and reports these residuals. The bridge
stores raw/scaled profiles, the common digest, daily residuals, and the
central-energy check; it is a deterministic realization of the fitted target
under the declared surrogate, not a second aggregate optimization.
For the job-to-network certificate, the indexed witness is mapped without a
second optimization. To avoid overloading the deadline endpoint $d_j$, let
$\mathcal J_{sk\delta}$ be the ledger partition whose source, class, and
declared destination labels are $s,k,\delta$; the evaluated witness uses the native
assignment $\delta=s_j$ and therefore has no outcome-dependent migration:
\begin{align}
 x^{\rm job}_{sk\delta\tau}&=\sum_{j\in\mathcal J_{sk\delta}:\,r_j\leq\tau<d_j}
 u_{j\tau},\\
 p^{\rm job}_{\delta\tau}&=P_\delta^{\rm fix}+\dt^{-1}X^{\rm job}_{\delta\tau},\qquad
 X^{\rm job}_{\delta\tau}=\sum_{s,k}x^{\rm job}_{sk\delta\tau}.
 \label{eq:jobnetworkcoupling}
\end{align}
Aggregate migration remains a matched-workload control; (\ref{eq:jobnetworkcoupling})
is the indexed replay profile. The audit reconstructs $X^{\rm job}$ from every
stored $u_{j\tau}$ and verifies equality before network valuation; no separately
optimized aggregate receives a network value. Workload deferral and
geographic assignment follow
\cite{cao2022flexibility,cao2024nonwire}, while rolling-state control follows
\cite{zhang2023receding}; waiting and
migration costs are $w_k(\tau-t)x_{skd\tau}$ and $m_{sd}x_{skd\tau}$, with
service-time substitution differing only by a fixed-arrival constant.

The submit-time and post-event ledgers have separate canonical digests. Each
submitted record is represented by the scheduler-only tuple; its Secure Hash Algorithm 256 (SHA-256) digest is
\begin{equation}
\begin{aligned}
 \ell_j^{\rm sub}={}&(\mathrm{id}_j,t_j^{\rm submit},\mathrm{timelimit}_j,
 g_j^{\rm req},\kappa_j,\sigma_j),\\
 \mathcal L_{\rm sub}={}&\operatorname{sort}\{\ell_j^{\rm sub}:j\in\mathcal J_{\rm sub}\},\\
 h_{\rm sub}={}&\operatorname{SHA\text{-}256}\!\left(\operatorname{serialize}(\mathcal L_{\rm sub})\right).
\end{aligned}
\label{eq:ledgerdigest}
\end{equation}
where $g_j^{\rm req}$, $\kappa_j$, and $\sigma_j$ are the requested-GPU, class,
and scheduler-state fields. The resulting canonical commitment is
(\ref{eq:ledgerdigest}). Start/end times, DCGM energy, and telemetry counts
enter a separate execution ledger whose join and energy equality are checked
only after commitment. Canonical ordering makes $h_{\rm sub}$ independent of
execution telemetry \cite{nist2015fips1804}. Measured power is reconciled with
the 118-MW nameplate fixed before the split.

\paragraph{Information and threat model.}
At decision time, $\mathcal I^{\rm dec}$ contains the submission ledger,
declared topology, pre-event prices, and configuration; the scoring set adds the
closed meter and execution ledger only after the contract is frozen. The threat
model covers post-commitment edits, replayed records, and attempts to introduce
future arrivals into a committed response. The canonical digest rejects an
altered scheduler tuple, the information-boundary certificate rejects future or
locked-test fields, and the settlement rule caps payment with the closed meter.
Scheduler and DCGM services remain trusted sources: a digest detects a changed
field, while upstream attestation is required to establish that a submitted
deadline, region, class, or energy field was semantically truthful. Complete-ledger
profiles are therefore ex-post information comparators; the slot-62 committed
ledger is the deployable contract protocol.

\subsection{Counterfactual, response, and false credit}

For event set $\mathcal E\subset\mathcal T$, $p$ denotes regional power in
MW and time sums use $\dt$ (MWh). The declaration-causal estimate is
$\widetilde p$; $p^0$ is the declaration-only strategic reference used for
mechanism isolation, $p^{1,\mathrm{sim}}$ is the independently solved
event-response schedule, and $p^{\rm obs}$ is the Massachusetts Institute of Technology (MIT) Data Center GPU Manager (DCGM) execution trace. The observed
trace is opened only after a profile is frozen; no utility event label exists,
so it is never a causal input. The mechanism-isolation response
$r=p^0-p^{1,\mathrm{sim}}$ is signed and retains rebound. For
$z\in\{\mathrm{obs},\mathrm{sim}\}$, submitted and true credit are
$\widehat r^{(z)}=[\widetilde p-p^{(z)}]_+$ and
$r^{(z)}=[p^0-p^{(z)}]_+$; false credit is
\begin{equation}
 F^{(z)}=\dt\sum_{d,t\in\mathcal E}
 [\widehat r^{(z)}_{dt}-r^{(z)}_{dt}]_+,
 \qquad z\in\{\mathrm{obs},\mathrm{sim}\}.
 \label{eq:falsecredit}
\end{equation}
The simulated source defines the validation epigraph; the observed source is
recomputed only in locked replay. False credit is computed from
(\ref{eq:falsecredit}) in absolute MW-slots, avoiding a ratio with a zero-credit
denominator. Before the event the operator freezes a contract baseline
$p^{\rm con}$. After meter closure, define the submitted gross reduction, the
contract-cap reduction, and the payable reduction by
\begin{equation}
\begin{aligned}
q^{\rm gross}_{dt}&=[\widetilde p_{dt}-p^{\rm obs}_{dt}]_+,\\
q^{\rm cap}_{dt}&=[p^{\rm con}_{dt}-p^{\rm obs}_{dt}]_+,\\
q^{\rm pay}_{dt}&=\min\{q^{\rm gross}_{dt},q^{\rm cap}_{dt}\}.
\end{aligned}
\label{eq:meter_cap}
\end{equation}
The payable response energy is
$Q^{\rm pay}_{\rm MWh}=\dt\sum_{d,t\in\mathcal E}q^{\rm pay}_{dt}$, and its
capacity cash component is $G^{\rm cap}=r^{\rm DR}Q^{\rm pay}_{\rm MWh}$.
Equation (\ref{eq:meter_cap}) uses only the submitted baseline, frozen contract,
and closed meter. Here $p^{\rm con}=p^{\rm safe}$ for the payment-band panel
and is the no-tariff committed-ledger profile for the decision-time panel;
$p^0$ remains an offline diagnostic. On locked event slots,
$\mathrm{nRMSE}=\sqrt{N_E^{-1}\sum_n(\widetilde p_n-p^0_n)^2}/
\max\{N_E^{-1}\sum_n p^0_n,\epsilon\}$. With
$M=\sum_n\min\{\widehat r_n,r_n\}$, precision, recall, and $F_1$ are
$M/\max\{\sum_n\widehat r_n,\epsilon\}$,
$M/\max\{\sum_n r_n,\epsilon\}$, and
$2PR/\max\{P+R,\epsilon\}$ over the $N_E$ region--slot samples.
The reported error metrics use mean absolute error (MAE) and root-mean-square error (RMSE); normalized RMSE (nRMSE) is the scale-free form above. The risk
simplex is the sole convex quadratic program (QP); all released workload
counterfactuals and payment caps use the linear programs stated above.

\subsection{Strategic reference scheduling}

Consider an arithmetic mean of $H=10$ reference days, event probability $q$,
and DR price $\pi^{\rm DR}$. If reference day $j$ adds event-window service
$\delta R_j$, the expected payment change is
\begin{equation}
 \Delta\Pi=\frac{q\pi^{\rm DR}}{H}\sum_{j=1}^{H}\delta R_j
 -\sum_{j=1}^{H}\Delta C_j .
 \label{eq:profit}
\end{equation}
Let $c'_j$ be the right derivative of minimum operating cost with respect to
the event-window service constraint on day $j$. Linear-program sensitivity
\cite{boyd2004convex} gives
\begin{equation}
\begin{aligned}
 \frac{q\pi^{\rm DR}}{H}>\min_j c'_j
 \quad\Longleftrightarrow\quad {}&
 \exists\ \text{a strictly profitable}\\[-2pt]
 &\text{local deviation}.
\end{aligned}
 \label{eq:threshold}
\end{equation}
At a degenerate optimum the current dual can be zero while the right marginal
cost is positive; an independent check evaluates $c'_j$ with $5\times10^{-4}$
and $10^{-3}$ MWh finite differences (agreement tolerance $10^{-6}$ USD/MWh). Equations (\ref{eq:profit})--(\ref{eq:threshold}) specialize
the historical-baseline premise in \cite{wang2022baseline}.

\subsection{Network value}

Let $M$ be the data-center-to-bus incidence matrix and
$p_t=p_t^{\rm native}+Mp_t^{\rm dc}$ the nodal demand vector. For
generator-to-bus map $C_g$, generator bounds, and power-transfer distribution
factor (PTDF) $H$ with line ratings $\bar f$, the
piecewise-linear security-constrained economic dispatch (SCED) value is
\begin{align}
 V(p)=\min_g\;&\sum_i C_i(g_i) \label{eq:sced}\\
 {\rm s.t.}\;&\mathbf1^\top g=\mathbf1^\top p,\quad
 \underline g\leq g\leq\bar g,\nonumber\\
 &-\bar f\leq H(C_gg-p)\leq\bar f.\nonumber
\end{align}
The incidence mapping aggregates demand at its connection bus; (\ref{eq:sced})
is the standard lossless DC representation \cite{stott2009dc,zimmerman2011matpower}.
The realized signed grid value is
\begin{equation}
 \Delta V=V(p^0)-V(p^1). \label{eq:value}
\end{equation}
Positive $\Delta V$ is a credit and negative $\Delta V$ a debit; no
positive-only truncation is applied.

\section{Proposed Methodology}\label{sec:method}

\subsection{Model-based framework and matched-information pre-estimation}

The state, power, and network equations define the feasible state before a
payment target is formed (Fig.~\ref{fig:framework}). Meter history and submitted
declarations enter the estimator in parallel; execution remains scoring-only.
Lagged meter/calendar baselines and declaration-view ledger comparators include
ridge, gradient boosting, and extremely randomized trees
\cite{hoerl1970ridge,friedman2001gradient}. Synthetic control and nonnegative
donor weighting are evaluated as a separate history-based class
\cite{geurts2006extratrees,abadie2010synthetic}.
The quantile output receives the same exact projection, with its penalty selected
by contiguous folds.

\begin{figure*}[t]
\centering
\includegraphics[width=0.82\textwidth]{figures/fig0_framework.pdf}
\caption{Verification-to-settlement framework. Validation labels select the
counterfactual risk contract; locked execution truth is used only for scoring.
The right branch computes aggregate N--1 value and exact participant shares.}
\label{fig:framework}
\end{figure*}

\subsection{Sparse cumulative workload-state verifier}

For statistical estimate $\widetilde p$ and each predeclared penalty $\rho_\ell$,
the first-stage program solves
\begin{equation}
 \min_{x,y,z}\ \mathcal C_{\rm w}(x)+\rho_\ell\sum_{d,t}z_{dt}
 \quad {\rm s.t.}\quad
 z_{dt}\geq\pm(p_{dt}(x)-\widetilde p_{dt}),
 \label{eq:projection}
\end{equation}
subject to (\ref{eq:cumulative})--(\ref{eq:power}). The absolute-value
epigraph is standard \cite{boyd2004convex}; its use for workload-ledger
verification is specific here. $\mathcal C_{\rm w}$ is workload cost, $C_i$ is
the generator-segment cost in (\ref{eq:sced}), and $V(p)$ is dispatch value, so
(\ref{eq:projection}) changes only the objective.
Six globally solved metadata schedules $p^{(\ell)}$ share identical arrivals
and constraints. A seventh causal Pareto projection is selected on validation
from the high-5-of-10 declaration target by a lexicographic rule over nRMSE,
false credit, under-credit, and $1-F_1$. The risk simplex contains these seven
causal profiles. A complete-ledger feasible-quantile projection is retained as
an external post-commitment comparator and is used only in the separate
payment hull. For a hull profile,
$\bar p=\sum_{\ell=1}^{L}\alpha_\ell p^{(\ell)}$, $\alpha\geq0$,
$\mathbf1^\top\alpha=1$, remains feasible because the common scheduling set is
convex.

The convex fit minimizes validation error plus total and daily-tail false-credit
exposure, measured by conditional value at risk (CVaR)
\cite{rockafellar2000cvar}. Let $\ell^\star$ minimize contiguous-fold nRMSE and
let $B_h$ be the absolute daily false-credit exposure of that single feasible
projection on validation day $h$. For a candidate simplex profile define
$F_h(\alpha)=\sum_{n\in h}f_n$, where $f_n$ is the positive-part false
credit at sample $n$. The risk fit uses absolute MW-slot exposure throughout;
there is no ratio with a potentially zero true-credit denominator. With
$m=\lceil(1-\gamma)H\rceil$ and reserve $\beta\in(0,1]$, solve
\begin{equation}
\begin{aligned}
\min_{\alpha,f,\nu,u}\quad
&\frac{\left\|\bar p-p^{0,\mathrm{decl}}\right\|_2^2}{N s_p^2}
 +\lambda\|\alpha\|_2^2
 +\omega_T\frac{\sum_{h=1}^{H}F_h}{S_T}\\[-2pt]
&\hspace{1.5em}+\omega_C\frac{\nu+m^{-1}\sum_{h=1}^{H}u_h}{S_C},
\end{aligned}
\label{eq:riskobjective}
\end{equation}
Equation (\ref{eq:riskobjective}) fixes both preferences inside the fit. Here
$s_p^2=\max\{N^{-1}\|p^{0,\mathrm{decl}}\|_2^2,\epsilon\}$,
$S_T=\max\{\sum_hB_h,\epsilon\}$, and
$S_C=\max\{m^{-1}\sum_{h\in\operatorname{tail}_\gamma}B_h,\epsilon\}$.
Single-axis ablations remove one term and its constraint; the locked joint fit
uses $\omega_T=1.00$, $\omega_C=0.50$, $\gamma=0.75$, and $\lambda=10^{-8}$.
The total-exposure term aggregates all event slots, while the CVaR term controls
the daily tail; both weights are fixed before locked scoring. The
budgets are
\begin{align}
 \sum_{h=1}^{H}F_h(\alpha)&\leq\beta\sum_{h=1}^{H}B_h,\label{eq:totalrisk}\\
 {\rm CVaR}_{0.75}\!\left(F_h(\alpha)\right)
 &\leq\beta{\rm CVaR}_{0.75}(B_h).\label{eq:cvar}
\end{align}
The positive part is represented without differentiating through a sort. For
each event sample $n$ and validation day $h$, introduce
$f_n\geq0$, a CVaR threshold $\nu\geq0$, and tail slacks $u_h\geq0$:
\begin{align}
 f_n&\geq p^\alpha_n-p^{(\mathrm{sim})}_n-r^{(\mathrm{sim})}_n,\\
F_h(\alpha)&=\sum_{n\in h}f_n,\\
 F_h(\alpha)-\nu-u_h&\leq0,\\
 \nu+\frac{1}{m}\sum_h u_h
 &\leq\beta\,\mathrm{CVaR}_{0.75}(B_h).
 \label{eq:riskepigraph}
\end{align}
The linear CVaR epigraph in (\ref{eq:riskepigraph}) uses the same $m$-day
denominator as the empirical statistic and avoids sorting. Here
$r^{(\mathrm{sim})}_n=[p^0_n-p^{(\mathrm{sim})}_n]_+$; the implementation uses
the equivalent threshold $p^{1,\mathrm{sim}}_n+r^{(\mathrm{sim})}_n$ and stores the
identity audit. To preserve the point-estimate metric, let
$q_\ell=N^{-1/2}\|p^{(\ell)}-p^{0,\mathrm{decl}}\|_2$ and
$q_\star$ denote the RMS error of the validation-selected causal Pareto
reference. The simplex additionally satisfies the linear sufficient condition
\begin{equation}
 \sum_{\ell=1}^{L}\alpha_\ell q_\ell
 \leq \sqrt{1+\varepsilon_{\rm acc}}\,q_\star,\qquad
 \varepsilon_{\rm acc}=0.10 .
 \label{eq:accuracycertificate}
\end{equation}
By the triangle inequality, (\ref{eq:accuracycertificate}) bounds the actual
simplex RMS error by $\sqrt{1+\varepsilon_{\rm acc}}q_\star$, equivalent to a
10\% validation-MSE budget. It therefore prevents a risk improvement from
being purchased with an unreported forecast degradation while retaining a
non-degenerate feasible simplex. The full problem is solved by the declared
CVXPY/Clarabel convex-QP path. Independent Karush--Kuhn--Tucker stationarity,
epigraph, simplex, accuracy, and primal residuals are recomputed after the
solve. Four contiguous folds select
$\beta$ only when each training contract is feasible. A separately reported
payment-contract projection targets $\bar p$ while imposing
\begin{equation}
p^{\rm safe}_{dt}\leq p^{(\ell^\star)}_{dt}. \label{eq:envelope}
\end{equation}
The pointwise cap in (\ref{eq:envelope}) is the upper side of the contractual
credit band; it is evaluated on every event slot before settlement and is not
substituted for the scientific risk profile.
\paragraph{Two-sided credit band.}
To exclude a zero-credit envelope, tie a lower band to the fitted target. Define
$p^{\rm floor}_{dt}=\min\{p^{(\ell^\star)}_{dt},\bar p_{dt}\}$, where $\bar p$ is
the pre-projection validation fit. The event-slot constraints are
\begin{equation}
 p^{\rm floor}_{dt}-\varepsilon
 \leq p^{\rm safe}_{dt}\leq p^{(\ell^\star)}_{dt},
 \qquad (d,t)\in\mathcal D\times\mathcal E,
 \label{eq:twosided}
\end{equation}
The release sets $\varepsilon=0$; the lower side is therefore the exact
pointwise minimum of the two independently formed profiles, and the selected
single projection proves nonemptiness. Equations
(\ref{eq:totalrisk})--(\ref{eq:cvar}) control aggregate and daily-tail exposure,
while the lower side in (\ref{eq:twosided}) regularizes under-credit and gives
no unseen-trace guarantee.
Realized overpayment is handled by (\ref{eq:meter_cap}).

\paragraph{Decision-time information boundary.}
A gate removes all post-gate arrivals before refitting the matched-information
target. The locked protocol uses slot 62, 30 minutes before the first event
slot, and a no-tariff LP freezes the contract baseline with a predeclared 3-MWh
minimum of participating event-window service. A second linear program (LP) applies the selected
positive DR price to the same masked ledger while carrying unexpired state.
Reserved future jobs cannot be credited, and the tariff plan is scored against
the closed meter through (\ref{eq:meter_cap}). The independent execution trace
is read only in the scoring pass.

To separate capacity planning from payment eligibility, the gate publishes a
causal reserve envelope. For candidate quantile $\eta$ and event day $d$,
\begin{equation}
 \widehat a^{\rm res}_{sk\tau}(d;\eta)=Q_{\eta}\!\left(\left\{a_{sk\tau}(h):h<d\right\}\right),
 \qquad \tau\geq t_{\rm gate},
\label{eq:causalreserve}
\end{equation}
Equation~(\ref{eq:causalreserve}) uses earlier days only. The predeclared
candidate set is $\{0.55,0.60,0.65,0.70\}$ with a 1.0-MWh/day validation
false-credit allowance; it selects $\eta=0.55$ before locked scoring. The
reserve is a planning envelope, and no reserved job enters payment before
commitment.



The cumulative state has
$O(|\mathcal S||\mathcal K||\mathcal D||\mathcal T|)$ nonzeros, versus quadratic
explicit prefixes. Equation~(\ref{eq:rollingterminal}) retains every observed
future arrival for the 512-slot deadline. Two lexicographic LPs (minimum L1,
then operating cost) embed the day trajectory and are solved globally with HiGHS
\cite{huangfu2018highs}. The resulting state-to-credit path is shown in
Figs.~\ref{fig:framework} and~\ref{fig:methoddetail}; execution and interval
audits are separate (Experiments~13--15).



\begin{figure}[t]
\centering
\includegraphics[width=\columnwidth]{figures/fig_method_detail.pdf}
\caption{Sparse cumulative verifier detail. Each admissible edge is a declared
release-to-deadline service opportunity; the global projection, risk fit, and
payment gate operate on the same feasible state and do not repair missing
deadlines or execution records.}
\label{fig:methoddetail}
\end{figure}

\subsection{Signed nodal and N--1 settlement}

Let $\lambda^0\in\partial V(p^0)$ be the nodal demand-value subgradient.
Signed linear settlement is $({\lambda^0})^\top(p^0-p^1)$, while
(\ref{eq:value}) is the realized signed system-value benchmark. Convexity
\cite{boyd2004convex} gives

\begin{equation}
 0\leq(\lambda^0)^\top(p^0-p^1)-\Delta V
 =V(p^1)-V(p^0)-(\lambda^0)^\top(p^1-p^0).
 \label{eq:bregman}
\end{equation}
Thus the linear rule upper-bounds exact value in the polyhedral market model;
equality requires a common supporting subgradient.

For market-compatible remuneration, let $\mathcal W$ contain the complete response cycle from anticipatory shift
through recovery. Its space--time remuneration is

\begin{equation}
 \Pi^{\rm ST}=\dt\sum_{t\in\mathcal W}\sum_d
 \lambda^0_{dt}(p^0_{dt}-p^1_{dt}).
 \label{eq:spacetime}
\end{equation}
This is the nodal virtual-link structure of
\cite{zhang2020virtuallinks,zhang2022remunerating}; summing site payments
reproduces (\ref{eq:spacetime}).

Space--time energy remuneration is one component of the contracted DR
product. Let $\mathcal D^{\rm part}$ be the predeclared participating sites,
$\mathcal E$ the event intervals, and $r^{\rm DR}$ the cleared service price.
The signed delivered capacity service and total operator value are
\begin{align}
 Q^{\rm cap}&=\dt\sum_{d\in\mathcal D^{\rm part}}\sum_{t\in\mathcal E}
 (p^\emptyset_{dt}-p^1_{dt}), \label{eq:capacityservice}\\
 G&=r^{\rm DR}Q^{\rm cap}+\Pi^{\rm ST}, \label{eq:operatorvalue}
\end{align}
where $p^\emptyset$ is the minimum-cost no-event schedule on the same window.
Equation~(\ref{eq:capacityservice}) defines the bilateral capacity product and
need not equal trace-anchored
$p^0$; $\Pi^{\rm ST}$ debits spatial transfer and recovery. This signed service
model follows demand-response mechanism design \cite{satchidanandan2023twostage,chen2021incentive}.

Let $x^1$ and $x^\emptyset$ induce $p^1$ and $p^\emptyset$. Define the
participant opportunity cost $C_{\mathrm{opp}}=\mathcal C_{\rm w}(x^1)-
\mathcal C_{\rm w}(x^\emptyset)$ and surplus $S=G-C_{\mathrm{opp}}$. With zero
disagreement utilities, the symmetric Nash contract \cite{nash1950bargaining} is

\begin{equation}
 z=\mathbb I\{S\geq0\},\quad
 P^{\rm B}=z\left(C_{\mathrm{opp}}+\frac{S}{2}\right),\quad
 u_{\rm dc}=u_{\rm op}=z\frac{S}{2}.
 \label{eq:nashcontract}
\end{equation}
Equation~(\ref{eq:operatorvalue}) combines the capacity payment and signed
space--time value. Equation~(\ref{eq:nashcontract}) uses nonactivation as the outside option,
keeps both utilities nonnegative, and is budget balanced. It is separate from
independent system operator (ISO) nodal remuneration and the exact dispatch-value benchmark.
For the executable full-day ledger, each interval first retains the signed
contract value $s_{dt}=\dt r^{\rm DR}q^{\rm pay}_{dt}+\Delta V_{dt}$, where
$q^{\rm pay}_{dt}$ is in MW, $r^{\rm DR}$ is in USD/MWh, $\dt$ is in hours,
and both terms are in USD. The reported economic
ledger is $S^{\rm signed}=\sum_{d,t}s_{dt}$; a cash presentation
$S^{\rm cash}=\sum_{d,t}\max\{s_{dt},0\}$ is stored separately. The cash floor
therefore leaves negative network-value debits visible in the signed cycle
total. The complete-cycle ledger $\mathcal W$ includes pre-event increases and
rebound, so workload conservation gives
$\sum_{d,t\in\mathcal W}(p^0_{dt}-p^1_{dt})\dt=0$ within solver tolerance.

Before payment comparison, the submitted, frozen-contract, and closed-meter
reductions are intersected as follows:
\begin{equation}
q^{\rm pay}_{dt}=\min\{q^{\rm obs}_{dt},q^{\rm cap}_{dt}\},\qquad
 G=r^{\rm DR}\dt\!\sum_{d,t\in\mathcal E}q^{\rm pay}_{dt}+\Pi^{\rm ST}.
\label{eq:settlementchain}
\end{equation}
Equation (\ref{eq:settlementchain}) is evaluated after the submitted profile,
contract, and closed meter are fixed; the second line uses
$q^{\rm pay}$ in the corresponding ledger. The observed-meter audit stores
the payment chain and its budget residual (below $10^{-8}$ USD). The
job-indexed audit stores the declaration replay
$q^{\rm replay}_{dt}=\min\{q^{\rm dec}_{dt},q^{\rm cap}_{dt}\}$ to certify
feasible network value. Exp30 records this role as declaration-only, and the
replay's signed/cash totals remain outside payable settlement until
$q^{\rm obs}$ is available.

Security-aware settlement augments (\ref{eq:sced}) with every finite
non-islanding single-line outage. If $L_{\ell k}$ is the LODF from outage $k$
to monitored line $\ell$, then
\begin{equation}
 -\bar f_\ell\leq
 (H_\ell+L_{\ell k}H_k)(C_gg-p)\leq\bar f_\ell,
 \quad \ell\neq k.
 \label{eq:n1}
\end{equation}
Equation (\ref{eq:n1}) is the standard preventive N--1 LODF formulation
\cite{tejada2018lodf}; every credible row uses native ratings. Islanding rows are
excluded because a connected single-reference PTDF cannot represent separate
islands, consistent with the DC/N--1 planning scope \cite{nerc2020tpl}.

\subsection{Scenario-robust N--1 payment certificate}

The pointwise envelope controls energy but not network payment because spatial
loads can activate different security constraints. Let $\mathcal P=\operatorname{conv}
\{p^{(1)},\ldots,p^{(L)}\}$ be the feasible-profile hull and
$p^{\rm cap}\in\mathcal P$ the validation-selected projection. Its penalty is
separate from the payment-target penalty; the feasible-quantile projection is an
external comparator. Let $\Xi_5$ be the 1st, 10th, 50th, 90th, and 99th
calibration-validation GPU-energy ratios from independent telemetry
\cite{samsi2021supercloud}. For facility profile $p^{\rm fac}_{dt}$, fixed demand
$P^{\rm fix}_d$, scale $s_{\rm dc}$, and factor $\xi$,

\begin{equation}
 p^{\rm dc,\xi}_{dt}=s_{\rm dc}\bigl(P^{\rm fix}_d+\xi[p^{\rm fac}_{dt}-P^{\rm fix}_d]\bigr),
 \qquad p^{\rm fac}_{dt}\geq P^{\rm fix}_d .
 \label{eq:fixedflexconversion}
\end{equation}
The scale uses the calibration-validation $q_{99}$ flexible peak, keeping the
fixed component invariant to $\xi$; raw endpoints remain activation scenarios.
The injection is $M p^{\rm dc,\xi}_t$, with superscripts $\alpha$ and ``cap''
applying (\ref{eq:fixedflexconversion}) to $p^\alpha$ and $p^{\rm cap}$. For
each settlement day, the \PaymentSafe{} defines
$B_\xi^{\rm cap}=\dt\sum_{t\in\mathcal E}
V_{N-1}(p_t^{\rm native}+M p_t^{\rm dc,\xi,cap})$ and solves
\begin{align}
\operatorname{lexmin}_{\alpha,e,\eta,\{g_t^\xi\}}\quad&
 \left(\sum_{d,t\in\mathcal E}e_{dt},-\eta\right)
 \label{eq:paymentcert}\\
\mathrm{s.t.}\quad&
 p^\alpha=\sum_{\ell=1}^{L}\alpha_\ell p^{(\ell)},\nonumber\\
&\alpha\geq0,\quad\mathbf1^\top\alpha=1,\quad 0\leq\eta\leq1,\nonumber\\
&e_{dt}\geq\pm(p^\alpha_{dt}-p^{\rm safe}_{dt}),\nonumber\\
&g_t^\xi\ \text{satisfies (\ref{eq:sced}) and (\ref{eq:n1})},\nonumber\\[-2pt]
&\hspace{2em}\text{with }p_t=p_t^{\rm native}+M p_t^{\rm dc,\xi,\alpha},\nonumber\\[-2pt]
&\hspace{7.5em}t\in\mathcal E,\quad \xi\in\Xi_5,\nonumber\\
&\dt\sum_{t\in\mathcal E}\sum_iC_i(g_{it}^\xi)\nonumber\\[-2pt]
&\hspace{2em}\leq(1-\eta)B_\xi^{\rm cap},
\quad \xi\in\Xi_5.\nonumber
\end{align}
The piecewise-linear epigraph is solved by two global LPs: trajectory error is
minimized first, then the common fractional margin $\eta$. Feasibility is the
N--1 upper bound and $p^{\rm cap}$ proves nonemptiness. Hence, for any
$\xi\in\Xi_5$ and realized load $p^{1,\xi}$, daily payment obeys
\begin{multline}
\sum_{t\in\mathcal E}\dt
[V_{N-1}(p_t^{\rm native}+M p_t^{\rm dc,\xi,\alpha})
-V_{N-1}(p^{1,\xi}_t)]\\
\leq\sum_{t\in\mathcal E}\dt
[V_{N-1}(p_t^{\rm native}+M p_t^{\rm dc,\xi,cap})
-V_{N-1}(p^{1,\xi}_t)].
\label{eq:paymentdominance}
\end{multline}
The realized cost cancels, so this is a samplewise certificate without execution
truth. Every finite non-islanding contingency is imposed using standard SCED/LP
primitives \cite{zimmerman2011matpower,boyd2004convex}. The hull contains six
first-stage projections and the matched feasible-quantile comparator. \RiskSafe{}
is the validation-fitted simplex point under the two risk budgets; the separately
stored payment contract targets it under the pointwise cap and is excluded from
risk fitting.

The finite cap also has a set-valued certificate. Let
$p^{\rm cap}=p^{(\ell^\star)}$, $p^{\rm pay}$ be the separately selected
payment target, and $p^{\rm cert}=p^\alpha$. For
$\Delta V_\xi(p)=V_{N-1}(p,\xi)-V_{N-1}(p^{1,\xi})$ and
$p_\theta=(1-\theta)p^{\rm cap}+\theta p^{\rm cert}$, the interval is
 \begin{equation}
  \mathcal I_\xi=\left[
  \min_{0\leq\theta\leq1}\Delta V_\xi(p_\theta),\;
  \max\{\Delta V_\xi(p^{\rm cap}),\Delta V_\xi(p^{\rm cert})\}\right].
  \label{eq:paymentinterval}
 \end{equation}
The lower endpoint in (\ref{eq:paymentinterval}) is one joint N--1 SCED LP with
a shared segment parameter; convexity gives the endpoint maximum and the cap
belongs to $\mathcal I_\xi$.
Validation-frozen profiles and the realized counterfactual determine both
endpoints; the interval width is uncertainty cost, while (\ref{eq:paymentcert})
remains the finite-scenario certificate.

\subsection{Participant allocation}

For participants $N$, let $v(S)$ be signed dispatch value when coalition $S$
moves from baseline to actual demand. The exact allocation is
\begin{equation}
 \phi_i(v)=\sum_{S\subseteq N\setminus\{i\}}
 \frac{|S|!(|N|-|S|-1)!}{|N|!}
 [v(S\cup\{i\})-v(S)].
 \label{eq:shapley}
\end{equation}
This Shapley specialization in (\ref{eq:shapley}) uses signed grid value
\cite{shapley1953value};
coalition-size grouping gives $\sum_i\phi_i=v(N)$ when $v(\varnothing)=0$.
All 16 four-site and 256 eight-portfolio coalitions are enumerated exactly.
For larger contracts, exchangeable site slices retain binomial multiplicities and
reduce evaluation to $\prod_g(m_g+1)$ count states (tested through 20 participants).

\subsection{Analytical guarantees}

These certificates delimit the logical scope. Under the interior and
differentiability assumptions of (\ref{eq:profit}), the threshold in
(\ref{eq:threshold}) is exact, as verified on the full 56-cell grid. The workload
set, simplex schedules, and two-sided envelope are feasible. For
any observed execution $c$, define credit from baseline $p$ as
$r(p;c)=[p-c]_+$ and the oracle baseline as $p^0$. The pointwise envelope in
(\ref{eq:twosided}) gives
\begin{align*}
[r(p^{\rm safe}_{dt};c_{dt})-r(p^0_{dt};c_{dt})]_+
&\leq[r(p^{(\ell^\star)}_{dt};c_{dt})-r(p^0_{dt};c_{dt})]_+,\\
[r(p^0_{dt};c_{dt})-r(p^{\rm safe}_{dt};c_{dt})]_+
&\leq[p^0_{dt}-p^{\rm floor}_{dt}]_++\varepsilon .
\end{align*}
The first inequality bounds excess credited response. For the second, $r(\cdot;c)$
is monotone and 1-Lipschitz, so the lower side of (\ref{eq:twosided}) bounds
missed credit by the oracle-to-floor gap and the declared band tolerance.
Summation gives the corresponding ledger bounds without locked outcomes.

Base-case and stacked N--1 values are convex piecewise-affine functions of nodal
demand; (\ref{eq:bregman}) gives the linear-value bound. Each dispatch embedded
in (\ref{eq:paymentcert}) is feasible, so (\ref{eq:paymentdominance}) holds for
the five frozen conversion scenarios $\Xi_5$. In the continuous segment audit, the realized counterfactual is fixed while
the affine segment between the two frozen contract profiles varies; a joint LP
gives the exact segment minimum, and convexity places its maximum at an
endpoint. No endpoint assertion is made when conversion simultaneously scales
the realized counterfactual. The payment-hull and segment audits use independent
ten- and four-segment models.

The digest detects post-commit scheduler-tuple changes, while the execution join
audits ordering and energy conservation. An optimum of (\ref{eq:jobindexed})
gives declaration-only runtime, capacity, and nameplate feasibility without
aggregate rounding. With $B_{dj}=\mathbf1\{s_j=d\}$ and bus map $M$, the typed
residuals are
\begin{align*}
r^{\rm job}_j&=\sum_{\tau=r_j}^{d_j-1}u_{j\tau}-E_j^{\rm run},\\
r^{\rm agg}_{d\tau}&=X_{d\tau}-\sum_jB_{dj}u_{j\tau},\\
r^{\rm map}_{b\tau}&=p_{b\tau}-P^{\rm fix}_{b\tau}-\dt^{-1}\sum_dM_{bd}X_{d\tau}.
\end{align*}
The first residual uses the same $E_j^{\rm run}$ as (\ref{eq:jobindexed});
$\widehat E_j$ remains diagnostic. Residuals at most $\varepsilon$ identify the
witness before N--1 valuation; allocation checks $\sum_i\phi_i=v(N)$ separately.
Moreover, if the feasible SCED value is $L_V$-Lipschitz on the committed
polytope, the mapped-value error is bounded by
\begin{equation}
 |V(p^{\rm map})-V(p^{\rm job})|
 \leq L_V\|r^{\rm map}\|_1 .
 \label{eq:typedbound}
\end{equation}
Equation~(\ref{eq:typedbound}) makes the zero-residual certificate explicit: the network solver receives
the same nodal profile implied by the indexed jobs, rather than an independently
refitted aggregate trace.

\section{Experimental Design}\label{sec:design}

\subsection{Data, preprocessing, and information boundary}

BurstGPT provides 5,188,507 successful requests and token service
\cite{wang2024burstgpt}; MIT Supercloud provides scheduler submissions and
independent Data Center GPU Manager (DCGM) energy \cite{samsi2021supercloud}.
The population contains 75,326 valid submissions and 71,128 positive-energy
execution matches; 68,662 have valid intervals in the common 121-day tensor.
Submission time defines arrivals, while measured intervals are scoring truth.
The GPU-power model uses 44,391 training, 9,637 calibration-validation, and
41,154 locked jobs. Calibration-validation ratios define the five scenarios
$\Xi_5$ in (\ref{eq:paymentcert}); execution joins are opened only after the
declaration witness is frozen. The 13 positive scheduler/DCGM pairs without a
valid interval are calibration labels only and cannot enter replay or selection.
The declared runtime plus a precommitted 96-slot queue allowance defines the
contract window; population membership is not filtered by outcomes.

The traces do not identify facility geography. Four-region coupling is
therefore a trace-driven benchmark with feature-stratified round-robin
assignment. We enumerate all 24 region-to-bus permutations and two concentration
controls at three penetrations. The PGLib native-load multiplier is fixed at
0.75 by the intact and contingency feasibility rule; all scored cells are
feasible. A separate fixed-generator-bus panel uses the pre-registered 0.90
multiplier and remains conditional on that scenario. Payment-network scale is
calibrated on $q_{99}$ at 4\% peak penetration, and all raw endpoints are
N--1-feasible under this scale. The flexible-only capacity-safe factor is
reported as a planning diagnostic, not an endpoint clip. PGLib, MATPOWER, and
RTS-24 supply the network cases \cite{babaeinejadsarookolaee2021pglib,zimmerman2011matpower,grigg1999rts}.

\subsection{Baselines, parameters, and inference}

The 54-day comparison includes meter and statistical baselines, feasible
quantile projections, synthetic control, the selected single projection, and
the joint and single-axis risk fits. Baselines share the declaration ledger;
execution truth is withheld until scoring. Deferral windows are 0, 4, and 512
slots, with 118-MW flexible nameplate and a 6-MW fixed component per region.
Projection penalties and normalized risk preferences are selected on
contiguous validation folds. The joint profile must satisfy a predeclared 10\%
MSE non-inferiority bound and total/CVaR exposure budgets. A separate slot-62
gate commits 3 MWh. Payment uses an independent ten-segment N--1 evaluator;
the contractual certificate uses four segments. $F_1$ is the harmonic mean of
precision and recall. 5,000-replication moving three-day bootstrap intervals and all $2^{18}$
block-sign assignments use Holm correction \cite{kunsch1989bootstrap,holm1979multiple}.

Sparse-system size ranges from 5,760 variables and 13,424 nonzeros at 96 slots
to 36,480 and 85,488 at 608 slots, with solve times of 0.039--0.212 s. Release
runs clamp each numerical backend to one thread; worker pools use at most 12
threads, below the 20-core limit. Long panels checkpoint schema- and
profile-validated results. The AC model enforces active/reactive balance,
voltage and reactive-generation bounds, and apparent-power limits
\cite{zimmerman2011matpower}. For each outage, non-reference active generation
is fixed to its intact value while the reference generator, reactive outputs,
and voltages remain recourse; all feasible cells pass the stated voltage,
thermal, and fixed-plan checks.

The evaluation spans strategic-baseline incentives, signed nodal settlement,
N--1 and AC feasibility, independent meter replay, declaration-indexed job
scheduling, and provenance checks. All comparisons use the same locked days
unless a deterministic network or ledger rule applies. An independent audit
reconstructs the job-to-region aggregation, verifies common-witness digests and
source hashes, and checks that risk, payment, and outage artifacts retain their
separate information roles.

\section{Results}\label{sec:results}

\subsection{Incentive boundary and independent verification}

The strategic-baseline test solves the entire $8\times7$ DR-price/call-probability grid.
The dual threshold (\ref{eq:threshold}) agrees with the independently solved
strategic program in all 56 cells. At \$150/MWh and $q=0.35$, strategic
reference scheduling inflates the event baseline by 8.223 MWh. Of
55.452 MWh credited, 44.458 MWh is unsupported by physical reduction,
giving a false-credit ratio of 0.802. The profitable region follows the derived
marginal boundary and does not use a fitted classifier; the complete grid and
its phase diagram are retained in the experiment output.

Twelve external and internal reference methods are evaluated on the locked
54-day mechanism-isolation panel (Table~\ref{tab:base}). Relative to the
metadata projection at $\rho=10$, the joint verifier reduces nRMSE, false
credit, and under-credit and raises $F_1$. That row is not the single profile
used to anchor the risk fit: the validation-selected Causal Pareto reference
has locked nRMSE/false credit/under-credit/$F_1$ of
0.6914/1.895/17.024/0.699. Against this internal reference, the joint fit
lowers false credit but has higher nRMSE and under-credit and lower $F_1$.
The total-budget-only ablation (0.8917/0.6455/26.3723/0.5490) improves all four
reported metrics over the joint fit and also has lower false-credit CVaR
(2.222 versus 2.607 MWh/day). Thus the locked panel does not establish that the
tail constraint adds predictive or empirical tail-risk benefit. The joint
profile remains the numerical optimizer of the declared seven-profile convex
program and meets its two validation-frozen budgets; those facts do not imply
dominance over the ablations on locked outcomes.
The validation certificate records reference/fitted MSE of 188.932/200.575
against the frozen 207.825 bound, KKT residual $5.80\times10^{-10}$, primal
residual $3.74\times10^{-11}$, and satisfied exposure budgets. Complete-ledger
rows are post-event diagnostics; they do not enter fitting or reserve selection.
On the 54-day paired three-day block bootstrap (5,000 draws), joint minus
total-only mean false credit is +0.134 MWh/day (95\% interval 0.016--0.287);
the 75\%-CVaR difference is +0.394 MWh/day (95\% interval -0.049--0.923).
The first interval excludes zero and the second includes it. The locked
ablation therefore gives no evidence that the CVaR term reduces false credit
or tail risk.

\begin{table}[t]
\centering
\caption{Locked 54-day mechanism-isolation response results (daily means).
External statistical comparators and internal risk references are shown
together; ex-post rows are information-boundary diagnostics.}
\label{tab:base}
\footnotesize
\resizebox{\columnwidth}{!}{%
\begin{tabular}{lrrrr}
\toprule
Method & nRMSE & False MWh & Under-credit & $F_1$\\
\midrule
High-5-of-10 & 1.586 & 41.220 & 11.169 & 0.572\\
Ridge & 1.038 & 21.906 & 25.226 & 0.545\\
Gradient Boosting & 1.011 & 20.977 & 26.401 & 0.534\\
Extra Trees & 1.090 & 26.629 & 18.176 & 0.597\\
Metadata Gradient Boosting & 0.774 & 9.183 & 21.528 & 0.640\\
Ex-post Metadata Gradient Boosting & 0.750 & 9.171 & 20.817 & 0.643\\
Ex-post Quantile Gradient Boosting & 1.337 & 0.591 & 54.391 & 0.237\\
Synthetic Control & 0.782 & 9.394 & 21.678 & 0.676\\
Feasible Quantile Projection & 1.344 & 0.098 & 52.647 & 0.225\\
Single Feasible Projection & 0.943 & 2.470 & 26.812 & 0.463\\
Causal Pareto risk reference & 0.6914 & 1.895 & 17.024 & 0.699\\
Total-budget-only verifier & 0.8917 & 0.646 & 26.372 & 0.549\\
Tail-Risk Feasible Counterfactual & 0.9051 & 0.824 & 26.548 & 0.532\\
\textbf{Risk-Constrained Convex Verifier} & \textbf{0.9021} & \textbf{0.780} & \textbf{26.506} & \textbf{0.536}\\
\bottomrule
\end{tabular}%
}
\end{table}

\begin{figure*}[t]
\centering
\includegraphics[width=0.47\textwidth]{figures/fig_risk_tail_tradeoff.pdf}\hfill
\includegraphics[width=0.47\textwidth]{figures/fig29_target_tracking.pdf}
\caption{Result evidence for the two coupled certificates. Left: held-out
normalized and absolute false-credit tails for the joint, total-only, and
CVaR-only fits. Right: the locked-day aggregate contract target and the
  declaration-indexed executable realization, with its audited residual below;
the realization is the exact finite-start policy minimizer, while target
tracking is audited rather than assumed.}
\label{fig:result_certificates}
\end{figure*}

The independent trace-meter replay yields event-window MAE of
1.909/6.880/7.038 MW for the feasible-quantile comparator/\RiskSafe/single
projection over the same 54 days, separate from mechanism-isolation scoring.
On the identical ledger, rolling-horizon,
batch-flexibility, and coupled-service translations obtain
0.943/2.470/0.463, 1.357/0/0.183, and 1.376/0/0.111
(nRMSE/false-credit/$F_1$), respectively; these are structural translations,
not software reproductions \cite{cao2022flexibility,chen2021idccoupling}.

\subsection{Nodal value and spatial response}

The settlement analysis crosses thirteen counterfactuals with five settlement mechanisms;
Table~\ref{tab:settlement} reports the trace-anchored locked rows. Exact signed
nodal value has MAE \$151.719 versus \$360.580 for nodal gross response, with
workload errors \$455.865 and \$421.677. Signed netting reduces trace error by
\$208.190/day (Holm-adjusted $p=3.89\times10^{-4}$); uniform-to-nodal pricing
changes it by $-\$1.463$ ($p=0.0118$), and linear-to-exact value by \$0.671
($p=0.8026$). All 5,616 interval/counterfactual instances satisfy
(\ref{eq:bregman}), with maximum residual $2.03\times10^{-11}$ dollars.

\begin{table}[t]
\centering
\caption{Base-case and complete N--1 settlement errors (USD/day).}
\label{tab:settlement}
\scriptsize
\resizebox{\columnwidth}{!}{%
\begin{tabular}{llrrcc}
\toprule
Panel & Settlement & Trace MAE & Workload MAE & Trace median/overpay & Workload median/overpay\\
\midrule
Base case & Nodal gross & 360.580 & 421.677 & -- & --\\
Base case & Nodal exact net value & \textbf{151.719} & 455.865 & -- & --\\
N--1 & Base exact & 139.221 & 167.367 & 79.79/130.42 & 121.01/107.69\\
N--1 & N--1 linear & 2.208 & \textbf{150.461} & 0.60/1.91 & 113.43/20.38\\
N--1 & \textbf{N--1 exact} & \textbf{0.990} & 150.955 & 0.87/0.01 & 113.43/20.18\\
\bottomrule
\end{tabular}%
}
\end{table}

The predeclared largest gross-minus-net offset (day 87) exposes the
rebound omitted by positive-only local accounting.

\subsection{Cross-network and N--1 security evidence}

Across 6,480 four-network outcomes, exact value improves trace error in all
12 cells (mean \$1.913/day; eight Holm-adjusted tests below 0.05), with a
\$0.050/day workload change. Capacity/deadline and coalition residuals are
below $5.68\times10^{-14}$ dollars. The 37-outage RTS-24 replay enforces
(\ref{eq:n1}), with trace/workload MAE
\$139.221/\$2.208/\$0.990 and \$167.367/\$150.461/\$150.955 and maximum
loading 1.000 p.u.

The complete replay scope is summarized in Table~\ref{tab:security}.

The payment certificate in (\ref{eq:paymentdominance}) is tested on 54
programs with a validation-frozen seven-profile hull. Contiguous folds select $\rho=10$ for the
cap and $\rho=3$ for the payment target; all five scaling-factor certificates
violate their bounds by less than $3\times10^{-11}$ USD. The
PaymentSafe/single/feasible-quantile/RiskSafe q50/q99 MAEs are
57.357/322.032, 41.141/231.841, 55.759/313.313, and 46.979/264.299 USD/day;
on two held-out interior factors, RiskSafe gives 9.345/14.042 versus
PaymentSafe's 9.988/15.006. Target selection excludes these transfer rows.
The independent evaluator uses ten segments; the contractual certificate uses
four, both fixed before locked scoring.

The AC replay has 864 connected-base and 5,400 corrective
IEEE-9/14 cells, all feasible. A fixed-plan contingency panel adds 9,612 method-domain
admissible fixed-plan contingencies and 144 intact solves with zero voltage
violation. The spatial-transfer panel evaluates 26 assignments at three penetrations under
the pre-registered 0.75 native-load PGLib domain; risk-constrained, single,
and feasible-quantile MAE triples are \$80.496/\$162.565/\$272.482,
\$89.537/\$180.649/\$300.873, and \$93.084/\$187.742/\$311.952.
Calibration $R^2$ is 0.8956/0.8196. Event-window MAE is
\$26.937/\$34.166/\$8.659/\$5.325, and complete-cycle residual is
\$0.621/\$0.569/\$0.104/\$0.

\subsection{Independent trace and deployment-boundary audits}

The immutable scheduler/DCGM replay contains 71,128 jobs, has a residual below
$10^{-17}$ MWh, and carries no event tag. The slot-62 evaluation scores the
complete-ledger comparator (0.902/0.780/0.536) and the committed-ledger gate
(1.366/0.494/0.080). The gate masks post-gate arrivals and execution truth
before scoring and requires 3 MWh on 54 days; its score therefore measures
information-boundary performance rather than payment. The independent
trace-meter MAE for \RiskSafe is 6.880 MW. Under a separate tariff, replay
precision/recall/$F_1$ is 0.0210/0.0194/0.529; independent event replay reaches
$F_1=1.000$. Table~\ref{tab:security} records the evidence boundary.

\begin{table}[t]
\centering
\caption{Information-boundary and certificate coverage. Triplets are
nRMSE/false response/$F_1$; declaration checks precede locked scoring.}
\label{tab:security}
\tiny
\setlength{\tabcolsep}{1.5pt}
\renewcommand{\arraystretch}{0.86}
\begin{tabular}{@{}p{0.46\columnwidth}p{0.45\columnwidth}@{}}
\toprule
Panel & Cells and certificate\\
\midrule
Complete-ledger joint fit & 54 days; 0.902/0.780/0.536; mechanism-isolation scoring\\
Slot-62 committed ledger & 54 days; 1.366/0.494/0.080; event-time profile\\
Indexed witness/bridge & 75,326 jobs; 13,198,247 variables; 7,306,622 starts; residual $<10^{-11}$ MWh; nRMSE 0.03023/relative-L2 1.581\\
DC N--1 replay & 864 cells; 37 outages/cell; 1.0000 p.u.; signed/cash 184.664/283.470 USD\\
AC certificates & IEEE-9/14: 864+5,400; four-network fixed-plan: 9,612 outages+144 intact; all admissible cells certified\\
\bottomrule
\end{tabular}
\end{table}

\subsection{Indexed workload counterfactual}

The declaration-indexed model assigns every valid submission a runtime-complete
block. Exact enumeration covers 7,306,622 feasible starts and reproduces the
selected minimizer (Table~\ref{tab:security}). Nameplate/central entitlements are 1,972.115/0.570
MWh, active declaration energy is 655.753 MWh, and energy residuals are below
$10^{-11}$ MWh. Event reduction is 4.68225 MWh. The submit-time-scaled
\RiskSafe{} bridge in (\ref{eq:bridge_residual}) has total nRMSE 0.03023, flexible relative-L2 1.581, and
maximum residual 3.220 MW. Its 5,184-cell ledger retains 1,060 debit cells,
network value $-0.686$ USD, signed cycle 184.664 USD, and cash-floor 283.470
USD. The payment audit labels this replay declaration-only; payment still requires a
closed meter.

\section{Conclusion}\label{sec:conclusion}

This paper couples a declaration-indexed workload verifier with signed nodal
valuation. On the complete-ledger mechanism-isolation panel, the joint
total-plus-CVaR fit reaches nRMSE/false credit/$F_1$ of 0.902/0.780/0.536 and
26.506 MWh/day under-credit, and meets its validation-frozen exposure budgets.
The locked ablation shows that the total-only fit improves all four metrics and
false-credit CVaR; the joint fit also does not dominate its own Causal Pareto
reference. These results leave the tail term as an unresolved design choice.
Figure~\ref{fig:result_certificates} reports the locked risk comparison and
the risk-to-executable residual. At the slot-62 gate, the committed-ledger
profile reaches nRMSE 1.366 and $F_1=0.080$, so full-ledger results do not
establish event-time performance. The task-level witness covers 75,326
submissions and 5,184 RTS-24 cells; its exact-start audit reports total-load
nRMSE 0.03023 but flexible relative-$L_2$ 1.581. The N--1 results are
declaration-only network replays, and a closed meter is required for payment.

\let\savedthebibliography\thebibliography
\let\savedendthebibliography\endthebibliography
% Keep the reference list readable while preserving the ten-page manuscript limit.
\renewenvironment{thebibliography}[1]{\savedthebibliography{#1}\fontsize{7}{7.2}\selectfont\setlength{\itemsep}{-0.8pt}\setlength{\parsep}{0pt}\setlength{\topsep}{0pt}}{\savedendthebibliography}
\bibliographystyle{ieeetr}
\bibliography{references}

\end{document}

